\begin{table}[!t]
\centering\footnotesize\setlength\tabcolsep{3pt}\renewcommand{\arraystretch}{1.08}
\caption{The defect screen: the five recorded result-invalidating defect kinds behind the records that the article marks as flagged. A record is flagged when its study carries one of these defects \emph{and} the defect touches that record's own claim (or its applicability is uncertain); a defect concerning another record of the same study leaves this one unmarked. The screen is the review's own rule, fixed before the findings were assembled; it is not a risk-of-bias instrument, and the article counts every record, flagged or not, reporting in S6 what setting the flagged records aside changes. Per-record judgements are listed in supplement S5e.}
\label{tab:defects}
\begin{tabular}{@{}L{1.9cm}L{2.4cm}L{2.0cm}r@{}}
\toprule
Defect & What it is & Why it invalidates the quantity & Rec. \\
\midrule
resubstitution & a probe or dictionary is fitted and scored on the same data & the reported quantity is an optimistically biased estimate of held-out performance, by an amount the data do not reveal & 8 \\
leakage across split units & a slide, patient or site appears on both sides of the split & the score is a within-unit score reported as an across-unit one & 7 \\
circular estimation & the quantity is estimated from a dependence the design created & the estimate is a property of the pipeline, not of the representation & 6 \\
selection on the reported metric & the configuration is chosen on the number reported & the number is a maximum over choices, reported as a single draw & 2 \\
single-case fitting & the fitting block is one case & nothing separates the case from the property & 1 \\
\midrule
Total & records with a recorded defect of these kinds & 19 touched, 1 uncertain, 4 left in & 24 \\
\bottomrule
\end{tabular}
\end{table}
